\begingroup
\emergencystretch=3em
\section{Метод решения}
\subsection{Эрозия и наращивание}
Пусть \(A\) --- исходная матрица размера \(N\times M\), а
\(r=\lfloor W/2\rfloor\) --- радиус окна. Для позиции \((i,j)\)
множество допустимых соседей определяется как
\[
\Omega_{ij}=\{(x,y):0\leq x<N,\ 0\leq y<M,\ |x-i|\leq r,\ |y-j|\leq r\}.
\]
Один проход эрозии и наращивания вычисляется по формулам
\[
E^{(0)}=A,\qquad D^{(0)}=A,
\]
\[
E^{(s+1)}_{ij}=\min_{(x,y)\in\Omega_{ij}} E^{(s)}_{xy},\qquad
D^{(s+1)}_{ij}=\max_{(x,y)\in\Omega_{ij}} D^{(s)}_{xy}.
\]
После \(K\) проходов возвращаются \(E^{(K)}\) и \(D^{(K)}\).
При \(W=1\) каждый элемент является единственным элементом своего окна,
поэтому обе операции сохраняют матрицу при любом числе проходов.

\subsection{Распределение работы}
Матрица разбивается по строкам между \(p=\min(P,N)\) участниками.
Участнику с номером \(q\), \(0\leq q<p\), назначается полуинтервал
\[
\left[\left\lfloor\frac{Nq}{p}\right\rfloor,
      \left\lfloor\frac{N(q+1)}{p}\right\rfloor\right).
\]
Диапазоны не пересекаются, покрывают все строки и отличаются по длине
не более чем на одну строку. Участник с номером 0 выполняется в главном
потоке; для остальных создаются \(p-1\) дополнительных потоков.
Каждый участник обрабатывает свой диапазон сразу для двух независимых
фильтров. Отдельные группы потоков для минимума и максимума не создаются.

\subsection{Буферы и синхронизация}
Для каждого фильтра выделяются два буфера. На проходе \(s\) читается буфер
\(s\bmod2\), а результат записывается в буфер \(1-(s\bmod2)\).
После прохода каждый участник ожидает остальные на барьере. Поэтому
никто не начинает чтение следующего поколения, пока оно не вычислено
целиком. По окончании возвращается буфер с индексом \(K\bmod2\).

Гонки при записи элементов отсутствуют: разные участники записывают
разные строки. Входной буфер текущего прохода только читается.
Мьютекс не нужен для каждого элемента матрицы; он применяется для
согласованного старта и неделимого вывода диагностических строк.

\subsection{Безопасный запуск потоков}
Перед созданием дополнительных потоков главный поток блокирует
стартовый мьютекс. Созданные участники ожидают его освобождения.
Если все потоки созданы, блокировка снимается и запускаются вычисления.
Если создание одного из потоков завершилось ошибкой, устанавливается
признак отмены старта. После снятия блокировки уже созданные участники
возвращаются, не входя в барьер. Главный поток ожидает их завершения
и сообщает об ошибке. Это исключает ожидание на барьере участника,
который не был создан.

\subsection{Оценка сложности}
При прямом переборе окна трудоёмкость двух фильтров имеет порядок
\(O(KNMW^2)\). Если окно превышает размеры матрицы, фактическое число
посещаемых соседей ограничено её размерами. В идеальной модели
вычислительная часть делится между \(p\) потоками, но к ней добавляются
создание потоков, синхронизация и ожидание планировщика.

Основные буферы содержат четыре матрицы \texttt{double}, что требует
примерно \(4NM\) элементов; размер одного элемента равен \texttt{sizeof(double)} байт.
Также выделяются задачи и дескрипторы потоков: \(O(p)\).
В эту оценку не входят стеки системных потоков и служебная память библиотек.

\section{Описание программы}
\subsection{Организация файлов}
\begin{longtable}{|p{0.25\linewidth}|p{0.66\linewidth}|}
\hline
Файл & Назначение \\ \hline
\texttt{common.hpp} & Параметры запуска, структура матрицы, общие функции. \\ \hline
\texttt{filters.hpp} & Структура результата и интерфейс \texttt{run\_filters}. \\ \hline
\texttt{main.cpp} & Разбор аргументов, чтение, запуск обработки, запись результатов. \\ \hline
\texttt{common.cpp} & Преобразование целых чисел, диагностика, PID/TID и таймер. \\ \hline
\texttt{matrix.cpp} & Выделение памяти, копирование, проверка и файловый ввод-вывод. \\ \hline
\texttt{filters.cpp} & Вычислительные задачи, двойная буферизация, запуск через \texttt{mt}. \\ \hline
\texttt{multithread/mt.h} & Общий интерфейс создания и ожидания потоков. \\ \hline
\texttt{multithread/}\newline\texttt{mt\_posix.cpp} & Linux-реализация функций интерфейса \texttt{mt}. \\ \hline
\texttt{mt\_sync.h},\newline\texttt{mt\_sync\_posix.cpp} & Обёртки над мьютексом и барьером POSIX. \\ \hline
\end{longtable}

\subsection{Основные структуры}
\texttt{Options} хранит путь входного файла, префикс выходных файлов,
ограничение числа потоков, число проходов, размер окна и признак
подробного вывода. \texttt{Matrix} содержит размеры и непрерывный массив
\texttt{double}, управляемый \texttt{std::unique\_ptr}. Матрица перемещается,
но не копируется неявно; для явного копирования используется
\texttt{copy\_matrix}.

\texttt{Context} содержит буферы двух фильтров, параметры вычислений,
барьер, мьютексы и признак отмены старта. \texttt{Task} задаёт контекст,
номер участника и границы его строк. \texttt{Result} возвращает две
матрицы, время вычислений и фактическое число участников.

\subsection{Интерфейс библиотеки потоков}
Основные операции интерфейса:
\begin{itemize}
\item \texttt{mt::CreateThread} --- создание потока;
\item \texttt{mt::ThreadJoin} --- ожидание завершения;
\item \texttt{mt::DestroyThreadHandle} --- освобождение дескриптора;
\item функции \texttt{mt::MutexCreate}, \texttt{mt::MutexLock}, \texttt{mt::MutexUnlock} и \texttt{mt::MutexDestroy} --- работа с мьютексом;
\item функции \texttt{mt::BarrierCreate}, \texttt{mt::BarrierWait} и \texttt{mt::BarrierDestroy} --- работа с барьером.
\end{itemize}
Платформенные операции скрыты в Linux-части библиотеки.
Сигнатуры создания и ожидания соответствуют интерфейсу проекта
\href{https://github.com/Yadroff/Operation-Systems}{репозитория преподавателя}.
Текущая реализация лабораторной рассчитана на Linux/WSL;
наличие пространства имён \texttt{mt} само по себе не означает
готовность всей лабораторной к нативной Windows.

\subsection{Функции ОС и библиотек}
\begin{itemize}
\item \texttt{pthread\_create}, \texttt{pthread\_join} --- создание и ожидание
системных потоков внутри Linux-реализации \texttt{mt}.
\item \texttt{pthread\_mutex\_*}, \texttt{pthread\_barrier\_*} --- взаимное
исключение и синхронизация проходов через обёртки библиотеки.
\item \texttt{clock\_gettime(CLOCK\_MONOTONIC)} --- измерение интервала
вычислений; системное календарное время для этого не используется.
\item \texttt{getpid}, \texttt{syscall(SYS\_gettid)} --- идентификаторы
процесса и Linux-потока для диагностики.
\item \texttt{sigaction}, \texttt{sigemptyset}, \texttt{write} --- обработка
\texttt{SIGPIPE} и диагностический вывод с повтором при \texttt{EINTR}.
\item \texttt{getopt}, \texttt{strtol}, \texttt{strtod},
\texttt{fopen}, \texttt{fscanf}, \texttt{fprintf}, \texttt{fclose} ---
функции библиотек для аргументов и файлового ввода-вывода.
\end{itemize}
Функции \texttt{pthread\_*} являются библиотечным интерфейсом, а не
одноимёнными системными вызовами ядра. Создание потоков в трассировке
Linux видно как \texttt{clone} или \texttt{clone3}; ожидания при конкуренции
могут использовать \texttt{futex}.
Семантика описана в
\href{https://man7.org/linux/man-pages/man3/pthread_create.3.html}{pthread\_create(3)},
\href{https://man7.org/linux/man-pages/man3/pthread_join.3.html}{pthread\_join(3)} и
\href{https://man7.org/linux/man-pages/man3/pthread_barrier_wait.3p.html}{pthread\_barrier\_wait(3p)}.

\subsection{Параметры запуска}
\begin{longtable}{|p{0.15\linewidth}|p{0.76\linewidth}|}
\hline
Ключ & Назначение \\ \hline
\texttt{-i FILE} & Обязательный путь входной матрицы. \\ \hline
\texttt{-t P} & Максимальное число потоков, включая главный; по умолчанию 1. \\ \hline
\texttt{-k K} & Число проходов каждого фильтра; по умолчанию 1. \\ \hline
\texttt{-w W} & Нечётная положительная сторона окна; по умолчанию 3. \\ \hline
\texttt{-o PREFIX} & Префикс результатов; по умолчанию \texttt{result}. \\ \hline
\texttt{-v} & Вывод PID, TID и диапазонов строк участников. \\ \hline
\texttt{-h} & Краткая справка по запуску. \\ \hline
\end{longtable}
\clearpage
\endgroup
